\section{文献综述}

本研究课题位于三个相互独立但有所交叉的学术文献的交汇处：（1）DeFi借贷协议的架构、风险及其实证研究；（2）行为金融理论，特别是以前景理论作为理解决策行为的框架；（3）关于链上信用风险评估以及利用区块链数据进行金融结果预测建模的新兴文献。以下逐一梳理这些领域的核心贡献，并指出其促成当前研究的特定空白。

\subsection{DeFi借贷：架构、风险与经验证据}

DeFi借贷协议无疑是去中心化金融生态系统中资金最密集、经验审查最充分的组成部分。这些协议通过智能合约实现无需许可的加密资产借贷，由算法自动管理抵押率、利率及触发清算事件，全程不依赖传统金融中介~\cite{schar2021decentralized2}。其基本机制如下：借款人存入加密资产作为抵押品，并可按其价值的一定比例（由协议的贷款价值比LTV参数定义）借出另一资产。如果抵押品价格相对借款金额下跌——或借出资产价格上涨——头寸的健康因子（HF）开始恶化。当HF跌至清算阈值（通常为1.0）以下时，任何网络参与者均可触发清算，偿还部分债务并索取相应抵押品及清算奖励~\cite{perez2020liquidations, qin2021empirical}。

从经验角度研究这些机制的学术文献已有相当积累。\citet{perez2020liquidations}对跨DeFi借贷平台的清算事件进行了最早的系统分析之一，记录了清算频繁发生、涉及大额交易费用并可能在涉及多个协议时引致连锁风险的现象。\citet{qin2021empirical}进一步考察了由清算奖金生成的激励结构，并展示了这些激励如何在高市场波动期间引发清算者之间的"矿工费优先竞拍"。\citet{sun2022liquidity}分析了Aave中流动性风险，揭示了大额头寸可能在特定资产池内造成显著的利用压力，从而可能限制小额储户的提款能力。

最近，一批文献开始关注跨协议风险及DeFi借贷的系统性维度。\citet{tovanich2023contagion}研究了Compound V2中的金融传染效应，构建了协议资产负债表并证明了集中化借贷头寸如何在相互连接的流动性池中产生不可忽略的违约风险，并通过网络传播。\citet{sevim2026interoperability}将风险建模扩展至多链环境，证明跨链互操作性引入了单链模型无法捕捉的新型风险类别。\citet{oberholzer2026toward}提出了面向机构级DeFi的九维综合风险评估框架，涵盖对手方风险、市场风险、流动性风险、操作风险、治理风险、法律风险、监管风险、智能合约风险以及预言机风险。\citet{iftikhar2025automated}对Aave和Compound中自动化风险管理机制进行了跨链比较分析，揭示了这两个主导借贷协议在参数调整和清算执行方式上的重要差异。

对该领域文献的一项回顾性特征在于其聚焦于协议层面或系统层面。不论其研究清算级联、利率动态还是跨协议依赖关系，其分析单元通常是协议、流动性池或交互中的智能合约集。个体借款人——作为其行为可能随风险敞口系统变化的决策者——主要是作为协议服务需求方而非自身成为研究对象的角色进入分析视野。因此，深入考察借款人如何主动应对风险积累以及其行为响应是否携带信用相关信息，构成了一项既有协议中心文献尚未系统探索的补充分析视角。

\subsection{前景理论与风险下的金融决策}

前景理论（Prospect Theory, PT）自\citet{kahneman1979prospect}首次提出并经\citet{tversky1992advances}进一步扩展为累积前景理论以来，一直是行为经济学中理解决策在风险与不确定条件下的最重要框架之一。该理论识别出了若干从预期效用理论预测中得出的系统性偏离。第一，个体以相对于参照点而非最终财富水平的方式评价结果。第二，价值函数在收益域呈凹形（风险规避），在损失域呈凸形（损失厌恶），损失的负效用约为等价收益效用的两倍。第三，个体呈现递减敏感性，即结果的边际感知随着远离参照点而递减。第四，决策权重在客观概率上系统性地呈非线性：低概率事件被高估，而中高概率事件被低估。

上述理论预测在广泛的金融领域中已经得到了充分的经验支持。\citet{benartzi1995myopic}揭示了近视损失厌恶——损失厌恶与频繁组合评估的结合——可以解释权益溢价之谜，而这是金融经济学中最突出的反常之一。\citet{barberis2013thirty}通过综合性综述记录了前景理论已经成功应用于解释股票交易中的处置效应、首次公开募股定价、权益溢价及一系列其他金融市场现象。该理论已成为行为金融的基石，提供了一个能够合理解释那些在标准理性预期范式中难以解释的经验规律。

然而，将前景理论应用于DeFi借贷市场既有机遇也存在挑战。一方面，DeFi借贷为测试行为理论提供了一个近乎理想的实验室。协议定义并通过编码写死在清算阈值参数中的健康因子值——头寸在此值以下即触达清算——充当了一个对所有参与者而言外生给定的客观参照点。这与传统金融领域形成鲜明对比，在后者中，参照点通常无法观测、具有主观性并可能在个体间差异显著~\cite{barberis2013thirty}。更进一步，DeFi中借款人所有操作——包括追加抵押品、还款、增加借款额和提取抵押品——的完全可观测性，意味着行为对清算阈值接近度的响应可以在比传统信用市场不可能实现精细粒度上进行度量。

另一方面，将前景理论应用于该情境也存在显著挑战。DeFi借款人不一定将协议定义的清算阈值视作其心理参照点；他们可能锚定在初始入场价格、组合层面目标回报率或其他基准上。此外，避免清算的理性动机——即支付给清算者的清算奖金形式带来的非可忽略财务处罚——本身就为围绕清算阈值的主动头寸管理提供了强有力的非行为解释。将理性风险管理与行为偏差区分开来，因此需要在经验设计上格外谨慎，包括利用控制组、工具变量或利用协议定义清算阈值生成拟外生变异的断点回归方法。

\subsection{链上信用风险评估}

与本课题相关的第三类文献涉及利用链上数据进行信用风险评估和违约预测的学术研究。这一新兴领域的核心前提在于：区块链数据的透明性使得在不依赖传统信用评分所依据的中央信用局、银行账户信息或传统金融记录的情况下构建信用风险模型成为可能。

近年来，已有若干研究论证了此一路径的可行性。\citet{ghosh2024onchain}开发了一种DeFi中链上信用风险评分框架，利用基于钱包交易历史、借贷和还款模式以及交互频率训练的机器学习模型来预测违约概率。其研究结果表明，即使在没有链下身份信息的情况下，链上行为数据也蕴含着显著的信用相关信息。\citet{xu2021banks}追溯了借贷市场从传统银行到去中心化金融的演进历程，认为DeFi协议有潜力通过实现更细粒度和更透明的风险评估来提高资本效率。若干综述文章汇集了可从链上数据提取的DeFi风险因素范围，包括跨多个协议的交易图谱、持币模式和交互历史~\cite{werner2022sok, jiang2023defi}。

然而，现有的链上信用风险文献存在若干本课题试图解决的关键不足。第一，现有研究主流将信用风险视作静态分类问题：鉴于在时间点\(t\)观察到的特征集，预测借款人在时间点\(t+k\)是否即将违约。虽然此一路径已取得有益成果，但其无法捕捉借款人在风险积累过中所经历的交易\emph{动态性}。一个健康因子在数周内逐渐下降的借款人如果在此过程中系统性地追加抵押品并降低杠杆，其风险状况将与具有完全相同时点健康因子但经历同等下跌幅度而无任何调整行为的借款人存在根本差异——然而静态模型会赋予他们相同的风险分数。

第二，现有文献聚焦于相对宽泛的链上活跃度度量指标，如总交易次数、账户年龄、跨协议的交互多样性以及协议层面聚合指标，而非聚焦于头寸接近风险阈值时发生的特定行为过程。状态变量（某一时点头寸当前的"快照"）与过程变量（头寸如何演变至当前状态以及借款人在此演变期间如何响应）之间的理论区分，虽然在更广泛的金融经济学文献中关于交易行为信息含量的讨论中地位突出，却尚未被系统地纳入到经验信用风险模型之中。

\subsection{识别研究空白}

上文对文献的梳理揭示了在三大研究领域的交叉处存在一个明确的研究空白。DeFi借贷研究从协议的视角已对清算机制、传染效应和系统性风险进行了全面记录，但尚未系统地考察借款人行为作为一个信用信号源。以前景理论为"封装标志"的行为金融学，为理解个体如何应对损失和参照点提供了一个丰富的理论框架，但尚未系统地应用于DeFi借贷语境。链上信用风险文献证明了区块链数据足以支持违约预测模型，但却采用了不能捕捉借款人风险管理过程性的主流静态性和基于特征的路径。

这一空白，可以明确地表述如下：\emph{目前尚无任何系统经验研究，考察借款人在DeFi借贷市场头寸风险积累期间的主动调整行为，是否包含了超越常规链上状态变量所捕捉的风险信息，以及所观察到的行为模式是否与前景理论的预测相符}。

填补这一空白在学术和实践层面均具有重要意义。在学术层面，检验行为过程变量是否携带增量预测能力，将连接传统上两个独立的研究流派将有望相互启发的以协议为中心的DeFi风险分析传统和以个体为中心的行为金融传统。在实践层面，如果证明主动借贷行为蕴含着独立的风险信号，将对协议设计产生直接含义，包括在用行为信息补充现有风险参数建立基于行为的早期预警指标方面的潜力。
